% The Binding Invariant: Synthetic Grounding as Cohomological Gradient Flow
% Flyxion -- standardgalactic
% April 2026

\documentclass[11pt, letterpaper]{article}

\usepackage{amsmath, amssymb, amsthm}
\usepackage{geometry}
\usepackage{microtype}
\usepackage{hyperref}
\usepackage{titlesec}
\usepackage{setspace}
\usepackage{mathtools}
\usepackage{tikz-cd}
\usepackage{enumitem}
\usepackage{abstract}
\usepackage{booktabs}
\usepackage{array}

\geometry{margin=1.25in}
\onehalfspacing

\hypersetup{
  colorlinks=true,
  linkcolor=black,
  urlcolor=black,
  citecolor=black,
  pdftitle={The Binding Invariant: Synthetic Grounding as Cohomological Gradient Flow},
  pdfauthor={Flyxion},
}

\titleformat{\section}{\large\bfseries}{\thesection.}{0.75em}{}
\titleformat{\subsection}{\normalsize\bfseries}{\thesubsection.}{0.6em}{}
\titleformat{\subsubsection}{\normalsize\itshape}{\thesubsubsection.}{0.5em}{}

%% Theorem environments
\theoremstyle{plain}
\newtheorem{theorem}{Theorem}[section]
\newtheorem{lemma}[theorem]{Lemma}
\newtheorem{corollary}[theorem]{Corollary}
\newtheorem{proposition}[theorem]{Proposition}

\theoremstyle{definition}
\newtheorem{definition}[theorem]{Definition}
\newtheorem{example}[theorem]{Example}

\theoremstyle{remark}
\newtheorem{remark}[theorem]{Remark}
\newtheorem{interpretation}[theorem]{Interpretation}

%% Notation
\newcommand{\BI}{\mathcal{B}}
\newcommand{\CC}{\mathcal{C}}
\newcommand{\HH}{\mathcal{H}}
\newcommand{\FF}{\mathcal{F}}
\newcommand{\MM}{\mathfrak{M}}
\newcommand{\EE}{\mathcal{E}}
\newcommand{\OO}{\mathcal{O}}
\newcommand{\Adm}{\mathcal{A}}
\newcommand{\traj}{\mathsf{traj}}
\newcommand{\op}{\mathrm{op}}
\newcommand{\Map}{\mathrm{Map}}
\newcommand{\holim}{\operatorname{holim}}
\newcommand{\hocolim}{\operatorname{hocolim}}
\newcommand{\Obs}{\mathrm{Obs}}
\newcommand{\grad}{\nabla}
\newcommand{\Grp}{\infty\text{-}\mathbf{Grpd}}
\newcommand{\argmin}{\operatorname*{argmin}}
\newcommand{\PSh}{\mathbf{PSh}}
\newcommand{\Sem}{\mathbf{Sem}}
\newcommand{\dSt}{\mathbf{dSt}}
\newcommand{\Id}{\mathrm{Id}}

\begin{document}

%% Title
\begin{center}
  {\LARGE\bfseries The Binding Invariant}\\[1.2em]
  {\large Synthetic Grounding as Cohomological Gradient Flow}\\[0.5em]
  {\normalsize\itshape
    Constraint Persistence, Derived Ambiguity, Process-Algebraic Equivalence,\\
    and Neural Garbage Collection as Semantic Renormalization}\\[1.8em]
  {\normalsize Flyxion}\\[0.4em]
  {\small April 2026}\\[2em]
\end{center}

\begin{abstract}
\noindent
Every system that undergoes irreversible transformation faces the same question:
what, if anything, carries through? This essay introduces the \emph{binding invariant}
as the structured property that persists under constraint closure when a system
passes from one regime of possibility to another.

We formalize this notion using constraint sheaves over partially ordered time and
prove the Binding Invariant Theorem: a system possesses a binding invariant if and
only if its constraint history admits a coherent section in the topos of feasible
trajectories. This yields precise criteria for uniqueness, stability, and rupture.

We then extend the framework using derived algebraic geometry, modeling ambiguity
as higher homotopy in a derived stack of semantic states. In this setting, the
binding invariant emerges as a homotopy-stable minimum of a cohomological energy
functional, and synthetic grounding is interpreted as a gradient flow that minimizes
obstruction across scales.

Bridging to concurrency theory, we show that Linear Logic Concurrent Constraint
(LCC) programming provides the operational counterpart of this geometry:
constraint consumption corresponds to obstruction reduction, and observational
equivalence (bisimulation and barbed congruence) corresponds to equivalence classes
of descent trajectories.

We further formalize CLIO as a monoidal endofunctor on a category of semantic states
and prove two central results: the associated energy functional is a
bisimulation-invariant Lyapunov function, and LCC bisimulation corresponds to
homotopy equivalence of descent trajectories.

Finally, we show that Neural Garbage Collection (NGC) provides an empirical
realization of the framework: its learned memory eviction dynamics implement
constraint-driven descent, its uncertainty signal behaves as a Lyapunov function,
and its graph-based reasoning approximates a homotopy colimit. A finite worked
example with diagrammatic proof grounds the abstract theory.

Taken together, these results establish a unified view in which identity is not a
static property of states, but the Lyapunov-stable equivalence class of trajectories
that remain coherent under irreversible constraint accumulation.
\end{abstract}

\newpage
%% ===========================================================================
%%  PART I
%% ===========================================================================

\begin{center}{\large\bfseries Part I: Foundational Framework}\end{center}
\vspace{0.5em}

\section{Introduction}

A system changes. Something survives. But \emph{what} survives, and in what sense?

This is not the ancient Ship of Theseus puzzle, though it rhymes with it. That problem
asks whether numerical identity persists under material replacement. The question here
is structurally deeper: it asks what \emph{constrains} the set of successors that a
system can coherently be---not what the system \emph{is} after transformation, but what
\emph{binds} its trajectory, making certain transformations continuations and others
ruptures.

We call this the \emph{binding invariant}: the property that must be preserved for a
transformation to count as development rather than replacement. More precisely, it is
the information-theoretic closure of the system's constraint history that any future
state must be consistent with in order to inherit the system's trajectory.

This question arises across multiple domains. In concurrency theory, observational
equivalence replaces state identity as the primary notion of process sameness
\cite{haemmerle2011}. In thermodynamics, irreversibility constrains possible
evolutions and places asymmetric costs on state transitions \cite{landauer,bennett,lloyd}.
In modern machine learning, systems that jointly learn reasoning and forgetting force a
reconceptualization of what internal state \emph{is} \cite{ngc2026}. The binding
invariant names what is preserved across all of these domains simultaneously.

The binding invariant arises at the intersection of three lines of work in the RSVP
theoretical program. \textbf{The Compiled Self} established that personal identity is
not a continuous substrate but a fixed point of a constraint-preserving adjunction over
an irreversible history topos; the binding invariant names what that fixed point
\emph{is}. \textbf{The Constraint--Recurrence Principle} showed that systems return
not to prior states but to prior constraint configurations, structured by a recurrence
operator; the binding invariant is the kernel of that operator. \textbf{The Yarncrawler
Identifiability Theorem} demonstrated that world-state reconstruction from partial
observations is possible if and only if the observation sheaf admits a global section;
the binding invariant is the analogue of that section in the context of
self-reconstruction across time.

The essay is organized as follows. Sections~2--4 (Part~I) develop the sheaf-theoretic
foundations and prove the Binding Invariant Theorem. Sections~5--8 (Part~II) develop
the derived-geometric and variational interpretation. Sections~9--11 (Part~III) bridge
to Linear Concurrent Constraints. Sections~12--14 (Part~IV) formalize CLIO as a
monoidal endofunctor with fixed-point algebra and sheaf-theoretic semantics.
Sections~15--17 (Part~V) prove the Lyapunov and bisimulation theorems. Sections~18--21
(Part~VI) absorb NGC and provide the worked example. A bibliography closes the essay.

\section{Formal Setting}

\subsection{The Site of Constraint Histories}

Let $S$ be a system evolving through a partially ordered time set $(T, \leq)$.

\begin{definition}[Site of Constraint Histories]
The \emph{constraint site} of $S$ is $\mathcal{J}_S = (T, \tau_{\mathrm{Alex}})$
where $\tau_{\mathrm{Alex}}$ is the Alexandrov topology on $(T, \leq)$, with coverings
given by upward-directed families $\{U_i\}$ such that $U = \bigcup_i U_i$ and
$t \in U_i, t \leq t' \Rightarrow t' \in U_i$.
\end{definition}

\begin{definition}[Constraint Sheaf]
A \emph{constraint sheaf} for $S$ is a sheaf $\CC$ on $\mathcal{J}_S$ such that:
(i)~the stalk $\CC_t$ is the set of constraints active at $t$; (ii)~restriction maps
$\rho_{t't}: \CC_{t'} \to \CC_t$ are monotone for $t \leq t'$ (the
\emph{irreversibility condition}: constraint satisfaction is preserved under forward
evolution); and (iii)~the gluing axiom holds.
\end{definition}

\begin{definition}[Trajectory Sheaf]
A trajectory is a section of a sheaf $\FF$ on $\mathcal{J}_S$ satisfying:
\[
  \FF(U) \;=\; \bigl\{\, \sigma : U \to \mathrm{States} \;\mid\; \sigma(t) \models \CC_t\,\bigr\}.
\]
\end{definition}

\begin{definition}[Constraint History]
The \emph{constraint history} of $S$ up to time $t$ is
$\HH_t(S) = \varprojlim_{t' \leq t} \CC_{t'}$.
\end{definition}

$\HH_t(S)$ is the structural sediment of all prior constraint regimes---generally
larger than $\CC_t$ because it retains the record of past closures even when present
constraints do not directly reflect them. Constraint-based models of computation
originate in Concurrent Constraint Programming, where systems evolve by monotonic
accumulation of information \cite{saraswat}. Linear logic extensions introduce
controlled consumption of constraints, breaking monotonicity while preserving logical
structure \cite{saraswatlincoln,girard}.

\begin{definition}[Closure Event]
A \emph{closure event} at time $t$ is a pair $(t, \gamma)$ where
$\gamma: \CC_t \to \CC_{t^+}$ is injective, generates strictly more constraints, and
is irreversible in the sense that no constraint-compatible $\gamma^{-1}$ exists as a
valid restriction map.
\end{definition}

\subsection{Topos of Feasible Trajectories}

Let $\mathbf{Traj}(S)$ be the category of constraint-compatible trajectories of $S$
and constraint-preserving maps between them.

\begin{proposition}
When $\CC$ satisfies the gluing condition, $\mathbf{Traj}(S)$ has the structure of a
Grothendieck topos.
\end{proposition}

\begin{proof}
The presheaf category $\widehat{T}$ is a topos. The constraint sheaf condition
restricts to the full subcategory of sheaves, which is a reflective subcategory of
$\widehat{T}$ and hence a topos by standard results in topos theory \cite{johnstone}.
\end{proof}

The subobject classifier $\Omega$ of $\mathbf{Traj}(S)$ classifies
constraint-admissible sub-trajectories: the truth values for the proposition ``this
sub-trajectory is a legitimate continuation of $S$.'' The use of $\Omega$ aligns the
construction with standard topos-theoretic semantics, where truth values classify
admissible substructures \cite{johnstone,maclane}.

\section{The Binding Invariant}

\begin{definition}[Binding Invariant]
The \emph{binding invariant} of $S$ at time $t$ is
\[
  \BI_t(S) \;=\; \Gamma\!\bigl(\HH_t(S),\;\Omega_{\mathbf{Traj}(S)}\bigr),
\]
the global sections of the subobject classifier over the constraint history.
\end{definition}

$\BI_t(S)$ is not a single property but the structured space of coherent judgments
about what counts as $S$ continuing. An element $b \in \BI_t(S)$ may be understood
as a rule assigning to each admissible sub-trajectory $U \subseteq T$ a truth value
indicating whether $U$ extends to a coherent global trajectory compatible with
$\HH_t(S)$. Thus $\BI_t(S)$ \emph{classifies admissible continuations} rather than
describing states.

\begin{definition}[Binding Section]
A \emph{binding section} at $(t, \gamma)$ is a section
$b \in \Gamma(\HH_t(S), \CC)$ such that $b(t') \in \CC_{t'}$ for all $t' \leq t$
and $\{b(t')\}$ is coherent in the inverse system.
\end{definition}

\begin{theorem}[Binding Invariant Theorem]
\label{thm:bi}
Let $S$ have irreversible constraint sheaf $\CC$ and closure sequence
$E = \{(t_i, \gamma_i)\}_{i=1}^n$ with effective descent for trajectories. The
following are equivalent:
\begin{enumerate}[label=(\roman*), noitemsep]
  \item $\BI_{t_n}(S) \neq \varnothing$.
  \item $\mathbf{Traj}(S)$ admits a global section compatible with $\CC$ after the
    last closure event.
  \item $\HH_{t_n}(S)$ admits a binding section compatible with all events in $E$.
  \item $E$ has a coherent limit in $\mathbf{Pro}(\mathbf{Traj}(S))$.
\end{enumerate}
\end{theorem}

\begin{proof}
The argument relies on standard sheaf-theoretic gluing and limit constructions
\cite{maclane,riehl}. (i)$\Rightarrow$(ii): By Yoneda, a global section of $\Omega$
over $\HH_{t_n}(S)$ corresponds to a non-empty object in $\mathbf{Traj}(S)$, which is
a global trajectory section. (ii)$\Rightarrow$(iii): The constraint components
$b(t') = \traj(\sigma(t')) \cap \CC_{t'}$ of any global section $\sigma$ form a
binding section, compatible with all closure events via compatibility of $\sigma$ with
the restriction maps of $\CC$. (iii)$\Rightarrow$(iv): A binding section defines a
compatible family of morphisms in $\mathbf{Traj}(S)$; the pro-limit exists by the
universal property. (iv)$\Rightarrow$(i): The pro-limit's projection maps define an
element of $\varprojlim_i \mathbf{Traj}(S)_{t_i}$; by the sheaf condition on $\CC$
this yields a global section of $\Omega$.
\end{proof}

\begin{corollary}[Uniqueness Criterion]
$\BI_t(S)$ is a singleton iff $\mathbf{Traj}(S)$ is Boolean at $t$.
\end{corollary}

\begin{corollary}[Rupture Criterion]
\label{cor:rupture}
A transformation at $t$ is a \emph{rupture} (replacement, not continuation) iff
$\BI_t(S) = \varnothing$.
\end{corollary}

\begin{corollary}[Binding Stability]
\label{cor:stability}
If $\BI_t(S) \neq \varnothing$ and all subsequent closure events are conservative,
then $\BI_{t'}(S) \neq \varnothing$ for all $t' > t$.
\end{corollary}

\begin{remark}[Degrees of Rupture]
Rupture may occur gradually. A weakening of the binding invariant corresponds to loss
of coherence in portions of $\HH_t(S)$, reflected in the non-existence of sections
over certain subdomains. Total rupture is $\BI_t(S) = \varnothing$; partial rupture is
failure of gluing on a covering. Identity is therefore not an intrinsic property of
states, but a fixed point in the category of constraint-consistent trajectories.
\end{remark}

\subsection{Dynamical Interpretation}

The binding invariant admits a dynamical interpretation that bridges to Parts~II--VI.
Let $\CC : X \to X$ be an update operator on semantic states, and let $\EE$ be an
energy functional measuring constraint inconsistency. A trajectory $x_{t+1} = \CC(x_t)$
induces a descent process $\EE(x_{t+1}) \leq \EE(x_t)$. Then $\BI_t(S)$ corresponds
to the set of equivalence classes of trajectories converging to the same minimal energy
configuration. The binding invariant is the fixed-point structure of a
constraint-induced descent flow.

\section{Applications and Philosophical Consequences}

\textbf{Cognitive identity.} The binding invariant gives the Compiled Self thesis a
precise interpretation: the self is the structured space of coherent judgments about
what would count as \emph{this person} continuing. Identity is not a substance but a
section; it is not determined by present state alone; rupture is formally possible; and
amnesia is structural---severe disruption of the binding section, not merely data loss.
This perspective aligns with modern machine learning systems in which reasoning and
memory are jointly optimized, and internal state is actively pruned during inference
\cite{ngc2026}.

\textbf{Computation.} A computation's binding invariant is the record of what it has
committed to. Na\"ive checkpointing fails when side effects alter the external
constraint environment in ways not captured by the checkpoint. System migration (not
mere update) is required when an upgrade destroys the binding section. In process
calculi, equivalence is defined observationally: two processes are identical if no
environment can distinguish them \cite{milner,haemmerle2011}.

\textbf{Physical processes.} First-order phase transitions are closure events; absence
of a binding section across the transition implies the post-transition system is a
distinct entity. The thermodynamic arrow of time is the direction along which binding
invariants accumulate: entropy production is the physical signature of irreversible
closure \cite{landauer,bennett,lloyd,jacobson}.

\textbf{Philosophical summary.} The binding invariant supports structural persistence
(what continues is a coherence condition, not a substance), refutes presentism (the
past is constitutive, not merely historical), and formalizes the distinction between
development and rupture. This position is structurally aligned with categorical and
structuralist approaches to identity, where objects are determined by their relations
rather than intrinsic properties \cite{spivak,riehl}.

%% ===========================================================================
%%  PART II
%% ===========================================================================

\vspace{1.2em}
\begin{center}{\large\bfseries Part II: Synthetic Grounding as Cohomological Gradient Flow}\end{center}
\vspace{0.5em}

\noindent\textit{Central claim}: A self, or coherent world-state, is a homotopy-stable
minimum of a cohomological energy functional defined over a derived moduli space of
candidate interpretations. The binding invariant is the structure that remains fixed at
such minima.

\section{Ambiguity as Higher Homotopy in Derived Stacks}

\subsection{The Derived Stack of Semantic States}

Let $\Omega$ denote the Spherepop event-log base. We model the space of candidate
interpretations as a derived stack:
\[
  \MM \;:\; \Omega^{\op} \;\longrightarrow\; \Grp,
\]
assigning to each open event-history $U \subseteq \Omega$ an $\infty$-groupoid
$\MM(U)$ of candidate interpretations. A ``thought''---a globally coherent
interpretation---is a section:
\[
  x \;\in\; \Gamma(\Omega, \MM) \;\simeq\; \holim_{U \subseteq \Omega} \MM(U).
\]

Instead of assigning to each $U_i$ a set of local sections, we assign an
$\infty$-groupoid whose objects are candidate local thoughts, whose 1-morphisms are
identifications between candidate thoughts, whose 2-morphisms are homotopies between
identifications, and so on. Ambiguity is not an external imperfection but intrinsic
higher structure. A global thought is a point in the homotopy limit
$\Gamma(\Omega, \MM) \simeq \holim_{U_i \in \mathcal{U}} \MM(U_i)$.

\subsection{Mapping Spaces and Ambiguity}

Two conflicting interpretations $x_1, x_2 \in \MM(U)$ are classified by the mapping
space $\Map_{\MM(U)}(x_1, x_2)$:
\begin{itemize}[noitemsep]
  \item \emph{Empty}: hard contradiction.
  \item \emph{Nontrivial}: ambiguity---a path exists but is underdetermined.
  \item \emph{Contractible}: equivalence---no meaningful distinction survives.
\end{itemize}

Higher-order ambiguity corresponds to higher homotopy. When paths between $x_1$ and
$x_2$ are themselves ambiguous, we ascend to 2-morphisms; ambiguity about 2-morphisms
yields 3-morphisms, and so on. The system carries this tower forward without forcing
premature resolution. Two candidate global thoughts need not be strictly identical in
their overlap data; they may be related by a path in
$\Map_{\MM(U_i \cap U_j)}(x_i|_{U_i \cap U_j}, x_j|_{U_i \cap U_j})$.

\begin{definition}[Derived Binding Invariant]
The \emph{derived binding invariant} of $S$ is
\[
  \mathbf{R}\BI(S) \;=\; \Map_{\dSt}(\HH_t(S),\;\MM).
\]
The classical $\BI_t(S)$ is recovered as $\pi_0(\mathbf{R}\BI(S))$.
\end{definition}

\subsection{The Self as $n$-Stack}

The binding invariant in the derived setting is the homotopy limit of the system's
homotopy-coherent diagram of candidate interpretations:
\[
  \BI_\infty(S) \;=\; \holim\bigl(\text{homotopy-coherent diagram of candidate sections}\bigr).
\]

\begin{theorem}[Derived Binding Invariant Theorem]
$\mathbf{R}\BI(S) \neq \varnothing$ iff the homotopy-coherent diagram of candidate
world-states over $\Omega$ has a non-trivial homotopy limit. When it exists,
$\BI_\infty(S)$ is an $n$-stack---a derived geometric object that persists despite
lower-dimensional fluctuations.
\end{theorem}

A point $x \in \MM$ has a tangent complex $T_x\MM$ encoding infinitesimal
deformations, while its higher cohomology encodes obstructions to extending local
deformations globally. A candidate thought is not merely present or absent: it has
nearby deformations, infinitesimal continuations, and obstruction classes preventing
extension into a stable global realization.

\section{The Cohomological Energy Functional}

\subsection{Obstruction Classes}

Let $H^\bullet(\Omega, \MM)$ be the cohomology of $\MM$ over $\Omega$. The degrees
carry distinct interpretations: $H^0$ recovers global sections; $H^1$ classifies local
inconsistencies (gluing failures, contradictions between adjacent event windows); $H^k$
for $k \geq 2$ classifies $k$-fold coherence failures (higher-order structural
ambiguity). Let $\Obs^k(x) \in H^k(\Omega, \MM)$ denote the $k$-th obstruction class
of state $x$.

A local version over a cover $\{U_i\}$ provides operational form:
\[
  \EE_{\mathcal{U}}(x)
  =
  \sum_{i < j}\alpha_{ij}\,d(x_i|_{U_i \cap U_j}, x_j|_{U_i \cap U_j})^2
  +
  \sum_{i<j<k}\beta_{ijk}\|\omega_{ijk}\|^2
  +
  \mu S(x) + \nu C(x),
\]
where the first term measures pairwise mismatch on overlaps and the second measures
higher cocycle defect on triple overlaps.

\subsection{The Global Functional}

Let $S(x)$ be semantic dispersion (entropy: the degree to which $x$ spreads
probability mass over incompatible interpretations), $C(x)$ be substrate cost (memory,
metabolic, or computational), $\mathcal{A}(x)$ the set of pending unsatisfied
constraints, and $r(a,x)$ the residual cost of pending constraint $a$ under state $x$.

\begin{definition}[Constraint Energy Functional]
\label{def:energy}
With parameters $\lambda_k, \mu, \nu, \eta > 0$:
\[
  \EE(x)
  \;=\;
  \sum_{k \geq 1} \lambda_k\,\bigl\|\Obs^k(x)\bigr\|^2
  \;+\;
  \mu\, S(x)
  \;+\;
  \nu\, C(x)
  \;+\;
  \eta \sum_{a \in \mathcal{A}(x)} r(a, x).
\]
\end{definition}

The four terms penalize distinct failure modes: failure of constraint closure at each
scale; semantic indeterminacy; unsustainable resource use; and residual unfulfilled
constraints (the LCC ``ask'' term, which drives the connection to process algebra in
Part~III). Low energy corresponds to vanishing obstruction, concentrated probability,
efficient representation, and satisfied constraints.

\section{Synthetic Grounding as Gradient Flow}

\subsection{CLIO Dynamics as Descent}

In the CLIO framework the system evolves via $x_{t+1} = \CC(x_t)$, where each step
restricts (prunes incompatible sections), infers (extends by constraint propagation),
and stabilizes. Geometrically: project $\to$ flow $\to$ stabilize.

\begin{definition}[Synthetic Grounding]
\emph{Synthetic grounding} is the stochastic gradient flow:
\[
  \frac{dx}{dt} \;=\; -\grad_x\,\EE(x),
  \qquad
  \mathbb{E}\bigl[\EE(x_{t+1})\bigr] \;\leq\; \mathbb{E}\bigl[\EE(x_t)\bigr],
\]
with strict decrease away from fixed points.
\end{definition}

Synthetic grounding is \emph{descent on cohomological inconsistency under resource
constraint}: the system prunes incompatible sections (reduces $\|\Obs^k\|^2$), merges
compatible ones (reduces $S(x)$), compresses redundancy (reduces $C(x)$), and
satisfies outstanding constraints (reduces the residual ask-cost term).

\subsection{Binding Invariants as Critical Points}

\begin{theorem}[Binding Invariants as Critical Points]
\label{thm:critical}
The critical points of $\partial_t x = -\grad\EE(x)$ are precisely the binding
invariants: states $x^*$ such that $[x^*] \in \BI_t(S)$.
\end{theorem}

\begin{proof}
At a critical point, $\grad\EE(x^*) = 0$. The gradient of the obstruction terms
vanishes iff $\Obs^k(x^*) = 0$ for all $k \geq 1$---$x^*$ is a globally coherent
section. The gradient of $\mu S(x)$ vanishes at entropy-minimum; $\nu C(x)$ at the
resource-efficient point; and $\eta\sum_a r(a,x)$ when all constraints are satisfied.
Together these are equivalent to $x^*$ being a global section of $\MM$ compatible with
the full event-history---a binding invariant.
\end{proof}

A binding invariant is:
\[
  x^* \;\in\; \argmin_{x \in \MM}\,\EE(x)
  \quad\text{such that}\quad x^* \simeq \CC(x^*).
\]
The self is a \emph{cohomologically minimal, homotopy-stable fixed point of CLIO
dynamics}.

\begin{theorem}[Convergence of Synthetic Grounding]
\label{thm:convergence}
If $\EE$ is convex on $H^\bullet(\Omega, \MM)$ and the event-log satisfies bounded
total variation, then $x(t) \to x^* \in \BI_t(S)$ as $t \to \infty$, uniquely when
$\EE$ is strictly convex.
\end{theorem}

\begin{remark}[Non-convex case and basin structure]
When $\EE$ is non-convex---as is typical for cognitive and social systems---the flow
may converge to any local minimum. Different constraint histories $\HH_t(S)$ define
different energy landscapes and hence different basins of attraction. CLIO is
\emph{locally contractive} within a basin:
$d(\CC(x), \CC(y)) \leq \kappa\, d(x,y)$ for $\kappa < 1$ when $x, y$ lie in the
same basin. Return Under Constraint is therefore guaranteed within a basin, not
globally. This explains identity fragmentation, allows multiple stable selves, and
aligns with LCC non-determinism.
\end{remark}

\section{Return Under Constraint as Geometric Stability}

\begin{proposition}
Every binding invariant $x^* \in \BI_t(S)$ satisfies the return property:
$\mathcal{R}_T(x^*) = x^*$ for all $T \geq 0$, where $\mathcal{R}_T = \Phi_T$ is the
time-$T$ map of the gradient flow.
\end{proposition}

\begin{theorem}[Geometric Persistence Theorem]
\label{thm:persist}
If $S'$ is obtained from $S$ by a conservative closure event $\tau$, then:
(i)~$S'$ has a binding invariant; (ii)~it is homotopic to that of $S$; (iii)~the
gradient flows of $\EE$ on $S$ and $S'$ are conjugate via an isomorphism induced by $\tau$.
\end{theorem}

\begin{theorem}[The Self as Cohomological Fixed Point]
\label{thm:self}
Under the convergence conditions of Theorem~\ref{thm:convergence}:
(i)~the self is $x^* = \lim_{t \to \infty} \Phi_t(x_0)$ for $x_0$ in the basin of
attraction; (ii)~$x^*$ is the fixed point of the return operator; (iii)~$x^*$
persists under $\tau$ iff $\tau$ is conservative; (iv)~rupture occurs iff $\tau$ moves
$x_0$ outside the basin of attraction.
\end{theorem}

The self compiles itself continuously: each incoming event perturbs $x$, and synthetic
grounding returns it to $x^*$. Identity is \emph{dynamical equilibrium}, not static
persistence.

\begin{center}
\begin{tabular}{@{}ll@{}}
\toprule
Cognitive / epistemic concept & Derived-geometric realization \\
\midrule
Ambiguity & Nontrivial homotopy in $\MM(U)$ \\
Contradiction & Empty mapping space \\
Reasoning & Descent on $\EE$ \\
Identity & Stable minimum of $\EE$ \\
Memory & Retained structure along gradient flow \\
Binding & Stability under recursive projection \\
Rupture & Exit from basin of attraction \\
\bottomrule
\end{tabular}
\end{center}

%% ===========================================================================
%%  PART III
%% ===========================================================================

\vspace{1.2em}
\begin{center}{\large\bfseries Part III: Process-Algebraic Correspondences}\end{center}
\vspace{0.5em}

\section{Linear Concurrent Constraints}

\subsection{Background}

Linear Concurrent Constraints (LCC), following Haemmerl\'e \cite{haemmerle2011}, is a
process algebra built on a constraint store that evolves by \emph{consumption}: the
$\mathsf{ask}$ operation reads a constraint from the store and destructs it, rewriting
the state dynamically. This breaks the monotonicity of classical constraint systems
\cite{saraswat} in precisely the way that CLIO dynamics breaks monotonicity via
eviction and projection. The key notions of observational equivalence in LCC are:
logical equivalence ($P^\dagger \vdash\dashv Q^\dagger$), may-testing, barbed
congruence ($\forall C,\; C[P] \sim C[Q]$), and labeled bisimulation. Haemmerl\'e
proves that labeled bisimilarity and barbed congruence coincide \cite{haemmerle2011}.

\subsection{Accessible Constraints as Binding Residue}

In LCC, the observable of a process $P$ is not its full state but
\[
  \OO_D(P) \;=\; \{\, c \mid P \;\Rightarrow\; P' \parallel c \,\}:
\]
the set of constraints that remain \emph{accessible} after evolution.

\begin{proposition}[Accessible Constraints = Binding Residue]
$\OO_D(P)$ is the operational shadow of the binding invariant: it corresponds to the
support of $\BI_t(P)$ in the constraint lattice.
\end{proposition}

The identification is structural, not analogical. The binding invariant is the
constraint-closed residue of the system's trajectory; $\OO_D(P)$ is the
constraint-closed residue of the process's execution. They are the same object viewed
from different levels of description.

\subsection{Constraint Consumption as Obstruction Projection}

\begin{proposition}
Constraint consumption in LCC corresponds to the projection
\[
  H^k(\Omega, \MM) \;\longrightarrow\;
  H^k(\Omega, \MM)\,/\,\langle\text{consumed constraints}\rangle.
\]
Each $\mathsf{ask}/\mathsf{consume}$ step modifies the cocycle data, reducing
$\|\Obs^k\|^2$ in expectation:
$\omega \mapsto \omega'$ with $\|\omega'\| \leq \|\omega\|$.
\end{proposition}

Constraint consumption is \emph{cohomological simplification}: it reduces obstruction
data, driving the gradient flow toward lower energy.

\begin{center}
\begin{tabular}{@{}ll@{}}
\toprule
LCC operation & Cohomological interpretation \\
\midrule
$\mathsf{tell}(c)$ & Add local section; may increase $H^1$ \\
$\mathsf{ask}(c)$ & Consume constraint; reduce $H^k$ obstruction \\
Parallel composition & Gluing of local sections \\
Commitment / choice & Irreversible closure event \\
Non-determinism & Branching of derived stack $\MM$ \\
\bottomrule
\end{tabular}
\end{center}

\section{Equivalences as Geometric Relations}

\subsection{Bisimulation as Gradient-Flow Equivalence}

\begin{definition}[Flow Bisimulation]
$P \approx_\EE Q$ iff for all flows $\Phi_t$, the trajectories $\Phi_t(P)$ and
$\Phi_t(Q)$ are homotopic in $H^\bullet(\Omega, \MM)$ at all times $t$.
\end{definition}

\subsection{Barbed Congruence as Descent Congruence}

\begin{definition}[Descent Congruence]
$P \cong_\EE Q$ iff for all perturbations $\delta$ (context embeddings),
$\CC(\delta(P)) \sim \CC(\delta(Q))$ under flow bisimulation.
\end{definition}

Barbed congruence is invariance under all gradient flows: no perturbation can separate
the binding invariants of $P$ and $Q$. This is equivalence under irreversible descent.

\section{CLIO-Observational Equivalence}

\begin{definition}[CLIO-Observational Equivalence]
\[
  x \;\sim_\CC\; y
  \;\;\Longleftrightarrow\;\;
  \lim_{t \to \infty}\CC^t(x) \;=\; \lim_{t \to \infty}\CC^t(y).
\]
Two states are equivalent if they converge to the same binding invariant.
\end{definition}

The refined asymptotic version avoids requiring a unique attractor:
\[
  x \sim_\CC y
  \;\Longleftrightarrow\;
  \lim_{t\to\infty} \mathrm{dist}\bigl(\CC^t(x), \CC^t(y)\bigr) = 0.
\]

\begin{definition}[Strong CLIO-Observational Equivalence]
$x \approx_\CC y$ iff $\forall t,\; \CC^t(x) \sim \CC^t(y)$. This is bisimulation
lifted to gradient flow: indistinguishability at every stage of descent.
\end{definition}

\begin{center}
\begin{tabular}{@{}lll@{}}
\toprule
LCC concept & Geometric realization & Formal object \\
\midrule
Constraints & Local sections & $\MM(U)$ \\
Accessible constraints & Binding residue & $\mathrm{supp}(\BI_t)$ \\
$\mathsf{ask}/\mathsf{consume}$ & Obstruction projection & $H^k \to H^k/\langle c\rangle$ \\
Bisimulation & Flow equivalence & $\Phi_t(P) \simeq \Phi_t(Q)$ \\
Barbed congruence & Invariance under irreversible descent & $\CC(\delta P) \sim \CC(\delta Q)$ \\
Logical equivalence & Derived stack equivalence & $\MM_P \simeq \MM_Q$ \\
May-testing & Shared basin of attraction & $\exists x \in \holim\MM$ reachable from both \\
\bottomrule
\end{tabular}
\end{center}

\begin{theorem}[Triple Characterization of the Binding Invariant]
\label{thm:char}
The binding invariant of $S$ admits three equivalent characterizations:
(i)~global sections of $\Omega$ over $\HH_t(S)$ (Part~I);
(ii)~homotopy-stable minima of $\EE$ (Part~II);
(iii)~the CLIO-observational equivalence class:
$\BI_t(S) \cong \{x \mid x \sim_\CC S\} / {\simeq}$ (Part~III).
\end{theorem}

%% ===========================================================================
%%  PART IV
%% ===========================================================================

\vspace{1.2em}
\begin{center}{\large\bfseries Part IV: CLIO as Monoidal Endofunctor}\end{center}
\vspace{0.5em}

\section{Semantic Category and Monoidal Structure}

Let $\Sem$ be a symmetric monoidal category whose objects are semantic states
$X \in \mathcal{X}$, whose morphisms $\phi: X \to Y$ are admissibility-preserving
transformations, and whose tensor product $\otimes$ represents compositional
aggregation of independent semantic components. We interpret $X = (\Phi, \mathbf{v}, S)$
as a structured RSVP field object in $\Sem$, where $\Phi$ is a scalar semantic
potential, $\mathbf{v}$ is a directional inferential flow, and $S$ is an entropy field
capturing unresolved structure.

\begin{definition}[CLIO Operator]
The \emph{CLIO operator} is a map $\CC : \mathcal{X} \to \mathcal{X}$:
\[
  \CC(X) \;=\; \mathbb{E}_{o \sim \pi_\theta(\cdot \mid E(X))}\bigl[\mathcal{F}(E(X), o)\bigr],
\]
composing: (i)~restriction $X \mapsto E(X)$; (ii)~inference $o \sim \pi_\theta$;
(iii)~state update $\mathcal{F}$.
\end{definition}

\begin{proposition}[Endofunctor Structure]
$\CC$ is an endofunctor on $\Sem$: for any morphism $\phi: X \to Y$,
$\CC(\phi): \CC(X) \to \CC(Y)$ is induced by functoriality of $\mathcal{F}$ and
invariance of $E$ over admissible maps.
\end{proposition}

\begin{proposition}[Lax Monoidality]
There exists a natural transformation
$\mu_{X,Y}: \CC(X) \otimes \CC(Y) \to \CC(X \otimes Y)$
making $\CC$ lax monoidal. Independent reasoning processes can be composed before or
after constraint closure, but the two are not strictly equivalent: interaction terms
arise through shared resource constraints.
\end{proposition}

\section{Fixed-Point Algebra of CLIO}

\begin{definition}[CLIO Algebra]
An algebra over $\CC$ is a pair $(X, \alpha)$ where $\alpha: \CC(X) \to X$.
\end{definition}

\begin{definition}[Fixed-Point Algebra]
A \emph{fixed-point algebra} satisfies $\alpha \circ \CC(\alpha) = \alpha$, and in the
special case $\alpha = \Id_X$, we obtain $\CC(X) = X$, or in the derived homotopy-correct
form, $X \simeq \CC(X)$.
\end{definition}

\begin{theorem}[Cognitive Fixed Point]
\label{thm:cogfixedpt}
Cognitive states correspond to fixed-point algebras of $\CC$: states invariant under
restriction (memory constraint), inference (generation), and update (state evolution).
A \emph{thought} is not a sequence but a fixed point of a recursive constraint-inference
operator.
\end{theorem}

\section{CLIO as Sheaf Endofunctor}

We lift $\CC$ from an endofunctor on semantic states to an endofunctor on presheaves:
\[
  \CC : \PSh(\Omega) \;\longrightarrow\; \PSh(\Omega),
\]
acting on sections by $\CC(s_U) = \mathbb{E}_o[\mathcal{F}(E(s_U), o)]$.

\begin{proposition}
$\CC$ preserves restriction maps up to stochastic equivalence, making it a weak
endofunctor on the category of presheaves.
\end{proposition}

\begin{definition}[Obstruction Class (Sheaf Form)]
Given a cover $\{U_i\}$ of $\Omega$, the obstruction to gluing local sections $\{s_i\}$
is $[\omega] \in \check{H}^1(\Omega, \FF)$.
\end{definition}

\begin{theorem}[Constraint Closure Criterion]
\label{thm:closurecrit}
A reasoning trajectory achieves constraint closure if and only if
$\check{H}^1(\Omega, \FF) = 0$.
Successful reasoning corresponds to the existence of a global section: all surviving
local computations are mutually compatible.
\end{theorem}

\begin{theorem}[Obstruction Descent]
\label{thm:obsdescent}
Under the entropy-reduction condition,
$\mathbb{E}[\|\check{H}^1(s^{(n+1)})\|] \leq \mathbb{E}[\|\check{H}^1(s^{(n)})\|]$.
Thus CLIO induces a descent process on cohomological obstruction.
\end{theorem}

\begin{proposition}[Functoriality under Admissible Maps]
If $\phi: \mathcal{X} \to \mathcal{X}$ preserves admissibility
($\phi(\Adm_B) \subseteq \Adm_B$), then
$\CC \circ \phi \approx \phi \circ \CC$
up to stochastic equivalence.
\end{proposition}

We obtain a unified picture of the KV-cache model: the cache defines a presheaf of
local semantic sections; reasoning generates local sections; eviction prunes sections
to reduce obstruction; CLIO iterates this as a functor; and cognition converges to
fixed points where obstruction vanishes.

%% ===========================================================================
%%  PART V
%% ===========================================================================

\vspace{1.2em}
\begin{center}{\large\bfseries Part V: Lyapunov Stability and Bisimulation Theorems}\end{center}
\vspace{0.5em}

\section{Setup and Assumptions}

Let $\mathcal{X}$ be a semantic state space equipped with CLIO update $\CC : X \to X$,
energy functional $\EE : X \to \mathbb{R}_{\geq 0}$ of the form in
Definition~\ref{def:energy}, and a realization map $F : \mathbf{LCC} \to X$ sending
LCC processes to semantic states. Let $\sim$ denote the observational equivalence
relation induced by LCC.

We operate under the following assumptions:
\begin{itemize}[noitemsep]
  \item[(A1)] \textbf{Observational invariance}: $P \sim Q \Rightarrow F(P)$ and $F(Q)$
    have the same accessible constraint image under $\OO_D$.
  \item[(A2)] \textbf{Bisimulation-compatibility of CLIO}: $x \sim y \Rightarrow
    \CC(x) \sim \CC(y)$.
  \item[(A3)] \textbf{Class-function property}: $x \sim y \Rightarrow \EE(x) = \EE(y)$.
  \item[(A4)] \textbf{Descent property}: $\EE(\CC(x)) \leq \EE(x)$ for all $x$,
    with strict inequality unless $x$ is in the stationary set.
  \item[(A5)] \textbf{Homotopy-coherent realization}: LCC transitions relevant for
    bisimulation yield paths $\gamma_{P,\alpha}: [0,1] \to X$ well-defined up to homotopy.
\end{itemize}

\begin{definition}[Admissible Fixed Point]
$X^* \in \Adm_B$ is an admissible fixed point if there exists $\pi_\theta$ such that
$\mathbb{E}[X_{t+1} \mid X_t = X^*] = X^*$ and $R(X^*)$ is locally maximal under
perturbations within $\Adm_B$.
\end{definition}

\begin{theorem}[Constrained Convergence]
\label{thm:cc}
Assume: (1)~$R: \mathcal{X} \to \mathbb{R}$ is bounded and Lipschitz on $\Adm_B$;
(2)~$E_t$ satisfies $\mathbb{E}[d(E_t(X), E_t(Y))] \leq \kappa\, d(X,Y)$ for $\kappa < 1$;
(3)~$\pi_\theta$ is updated via an unbiased policy gradient with diminishing step size.
Then $\{X_t\}$ converges in distribution to a set of admissible fixed points.
\end{theorem}

\begin{proof}
The contraction condition prevents unbounded growth of state variability. Boundedness
and Lipschitz continuity of $R$ ensure stability under restriction. Standard stochastic
approximation results imply convergence of $\pi_\theta$ to a stationary policy.
Together, these imply that the coupled process $(X_t, \pi_\theta)$ converges to
stationary points where both state and policy are invariant in expectation, corresponding
to admissible fixed points.
\end{proof}

\begin{proposition}[Entropy Reduction under Restriction]
\label{prop:entropy}
If $E_t$ removes components whose expected contribution to future reward is
non-positive, then $\mathbb{E}[S_{t+1}] \leq \mathbb{E}[S_t]$, with strict inequality
whenever non-contributing components are present.
\end{proposition}

\section{Theorem 1: Bisimulation-Invariant Lyapunov Function}

\begin{theorem}[CLIO Induces a Bisimulation-Invariant Lyapunov Function]
\label{thm:lyapunov}
Under assumptions A2--A4, $\EE$ is a Lyapunov function on the quotient space
$\mathcal{X}/{\sim}$, and therefore a bisimulation-invariant Lyapunov function for
CLIO.
\end{theorem}

\begin{proof}
\textbf{Step 1.} By A3, $\EE$ descends to a well-defined function
$\bar{\EE}: \mathcal{X}/{\sim} \to \mathbb{R}_{\geq 0}$, setting
$\bar{\EE}([x]) := \EE(x)$. This is well-defined because $x \sim y$ implies
$\EE(x) = \EE(y)$.

\textbf{Step 2.} By A2, $\CC$ induces a quotient map
$\bar{\CC}: \mathcal{X}/{\sim} \to \mathcal{X}/{\sim}$, setting
$\bar{\CC}([x]) := [\CC(x)]$, which is well-defined because $x \sim y$ implies
$\CC(x) \sim \CC(y)$.

\textbf{Step 3.} Apply A4:
$\bar{\EE}(\bar{\CC}([x])) = \EE(\CC(x)) \leq \EE(x) = \bar{\EE}([x])$.
So $\bar{\EE}$ is non-increasing along CLIO iterates on equivalence classes. If $[x]$
is not stationary, the decrease is strict, so $\bar{\EE}$ is a strict Lyapunov
function away from the stationary set.

Therefore $\EE$ is Lyapunov on $\mathcal{X}/{\sim}$ modulo bisimulation. \qed
\end{proof}

\begin{corollary}[Stationary Classes are Binding Invariants]
If $\bar{\CC}([x^*]) = [x^*]$ and $\bar{\EE}$ is minimized at $[x^*]$, then $[x^*]$
is a binding invariant in the strong sense: a stable observational class that no
admissible CLIO update can distinguish from itself.
\end{corollary}

\begin{interpretation}
Identity is not assigned at the level of raw presentation but at the level of
observational class. $\EE$ does not care which representative is chosen, and CLIO
dissipates it monotonically. What stabilizes is not a state snapshot but an equivalence
class under interaction---in Haemmerl\'e's language, a bisimulation class; in ours, a
binding invariant.
\end{interpretation}

\section{Theorem 2: Bisimulation as Homotopy Equivalence of Descent Trajectories}

\begin{theorem}[LCC Bisimulation = Homotopy Equivalence of Descent Trajectories]
\label{thm:bisim-homotopy}
Assume A1, A2, A5. Let $P, Q$ be LCC processes such that $P \approx Q$ (labeled
bisimilarity in Haemmerl\'e's sense). Then the CLIO descent trajectories issued from
$F(P)$ and $F(Q)$ are homotopy equivalent in the quotient dynamics: for every
bisimulation-matching transition sequence from $P$, there is a matching sequence from
$Q$ whose realized geometric trajectory is homotopic relative endpoints in $X$.
\end{theorem}

\begin{proof}
Labeled bisimulation means that if $P \xrightarrow{\alpha} P_1$, then there exists
$Q_1$ such that $Q \xRightarrow{\alpha} Q_1$ and $P_1 \approx Q_1$ (Haemmerl\'e
\cite{haemmerle2011}, Definition of DE-bisimulation). Apply $F$. By A5, the transition
$P \xrightarrow{\alpha} P_1$ yields a path $\gamma_{P,\alpha}: [0,1] \to X$ from
$F(P)$ to $F(P_1)$; the matched transition yields $\gamma_{Q,\alpha}$ from $F(Q)$ to
$F(Q_1)$. By A1 and bisimulation, $F(P) \sim F(Q)$ and $F(P_1) \sim F(Q_1)$.
Therefore in the quotient space $\mathcal{X}/{\sim}$, both paths represent the same
arrow $[F(P)] \to [F(P_1)]$ up to homotopy. Iterating over all matched transitions
yields a ladder diagram whose squares commute up to homotopy. By standard pasting, the
full descent trajectories are homotopy equivalent relative endpoints. \qed
\end{proof}

\begin{corollary}[Barbed Congruence and Homotopy]
Since Haemmerl\'e proves that labeled bisimilarity and barbed congruence coincide
\cite{haemmerle2011}, the homotopy classification above applies equally to barbed
congruence:
$P \cong Q \;\Longleftrightarrow\; F(P), F(Q)$ determine the same homotopy class of
descent trajectories. This is especially important because barbed congruence is the
irreversible, committed-choice notion---exactly matching the RSVP program's emphasis on
projection, path dependence, and return under constraint.
\end{corollary}

\begin{theorem}[Quotient Descent on Observational Classes]
\label{thm:quotient-descent}
Under A1--A5, CLIO defines a descent dynamics on the homotopy category of
observational equivalence classes, and the induced $\EE$ is a Lyapunov function on
that category.
\end{theorem}

\begin{proof}
From Theorem~\ref{thm:lyapunov}, $\EE$ descends to a Lyapunov function on
$\mathcal{X}/{\sim}$. From Theorem~\ref{thm:bisim-homotopy}, the morphisms generated
by CLIO trajectories in $\mathcal{X}/{\sim}$ are precisely homotopy classes of
operationally equivalent LCC evolutions. The quotient dynamics is simultaneously
operationally meaningful and geometrically coherent.
\end{proof}

\noindent\textbf{Final compression of Part~V.} Two claims are now fully established:
\begin{quote}
\itshape
(1) $\EE(\CC(x)) \leq \EE(x)$ is not merely an energy drop on raw states---it is an
energy drop on bisimulation classes. The Lyapunov object is identity-respecting.

(2) $P \approx Q$ does not merely mean two programs cannot be told apart by an
observer. Under realization into the derived semantic space, it means their descent
dynamics belong to the same homotopy class.
\end{quote}

%% ===========================================================================
%%  PART VI
%% ===========================================================================

\vspace{1.2em}
\begin{center}{\large\bfseries Part VI: Neural Garbage Collection as Constraint-Driven State Reduction}\end{center}
\vspace{0.5em}

\section{From Monotonic Context Growth to Adaptive State Restriction}

Standard autoregressive language models operate under a monotonic growth assumption:
the internal state $X_t$ expands strictly with each inference step:
\[
  X_{t+1} \;=\; \mathcal{F}(X_t, o_t),
  \qquad X_0 \subseteq X_1 \subseteq X_2 \subseteq \cdots
\]
This presupposes that all previously generated information remains equally available
and relevant.

Neural Garbage Collection (NGC) \cite{ngc2026} introduces a fundamental deviation. At
discrete intervals, the model applies a learned restriction operator
$E_t: X_t \to X_t'$, yielding:
\[
  X_{t+1} \;=\; \mathcal{F}(E_t(X_t), o_t),
\]
breaking monotonicity. The model jointly learns \emph{what to think} and \emph{what to
forget} under a single reward signal. The eviction decisions are not heuristic (like
recency or attention thresholds) but are learned via reinforcement learning using only
final task success. The most important sentence in the paper is essentially:
\begin{quote}\itshape
``If a model can learn to reason, why can't it learn to forget?''
\end{quote}
This is a conceptual inversion: previously, reasoning was learned and memory was
engineered; now reasoning and memory are co-learned.

\subsection{RSVP Field Representation}

We reinterpret the internal state as an RSVP field:
\[
  X_t \;=\; (\Phi_t, \mathbf{v}_t, S_t),
\]
where $\Phi_t$ is the scalar semantic potential, $\mathbf{v}_t$ is the directional
inferential flow, and $S_t$ is the unresolved entropy field. The KV cache is a
discrete support over a domain $\Omega_t$:
\[
  X_t \;\approx\; X|_{\Omega_t}.
\]
Eviction corresponds to selection of a subdomain $K_t \subset \Omega_t$, yielding a
restricted field:
\[
  X_t' \;=\; X|_{K_t},
  \qquad
  X_{t+1} \;=\; \mathcal{F}(X|_{K_t}, o_t).
\]
Reasoning proceeds not on the full accumulated field, but on a dynamically selected
support: inference is no longer transport on a fixed field history, but transport on a
\emph{self-pruned field manifold}.

\subsection{Variational Formulation}

Let $R(X_T)$ denote terminal task realizability and $C(K_t)$ retained computational
substrate. The joint optimization over reasoning and eviction is:
\[
  \max_{\{o_t, K_t\}}\; R(X_T)
  \quad\text{subject to}\quad
  \sum_t C(K_t) \;\leq\; B.
\]
Introducing a Lagrange multiplier $\lambda$, we obtain the unconstrained functional:
\[
  \mathcal{J} \;=\; R(X_T) \;-\; \lambda\sum_t C(K_t).
\]
In RSVP language, this is a constrained variational principle: reasoning and forgetting
are coupled controls minimizing future work while preserving task closure. This is
precisely the ``selection determines reality'' thesis: the realized trajectory is the
residue of a decimation process.

\subsection{Eviction as Entropic Projection}

\begin{definition}[Admissible Projection]
Eviction is interpreted as an approximate projection
\[
  E_t: X_t \;\mapsto\; \Pi_{\Adm}(X_t),
\]
where $\Adm$ denotes the admissible subspace of states preserving downstream
realizability under resource constraints.
\end{definition}

NGC implements \emph{entropic decimation}: components of the state that do not
contribute to future constraint closure are removed. The realized trajectory is not the
accumulation of all intermediate steps, but the residue of selective pruning.

\subsection{Sheaf-Theoretic Interpretation}

Let $\{U_i\}$ be a cover of the reasoning trajectory, where each $U_i$ corresponds to
a contiguous block of reasoning tokens. The KV cache encodes local sections
$s_i \in \FF(U_i)$. NGC performs a learned pruning $\{s_i\} \mapsto \{s_i\}_{i \in I_t}$
of sections that are likely to admit extension to a global section under future
constraints. Eviction functions as a coarse, learned filter on the gluing problem,
discarding sections unlikely to contribute to eventual consistency. This interpretation
is especially natural because NGC performs eviction at the level of contiguous token
blocks, which correspond to semantically coherent subcomputations occupying contiguous
spans.

More precisely, eviction is an \emph{approximate obstruction-sensitive restriction
functor}: the model is learning, however crudely, which local sections need not be
preserved for later consistency. The stronger formulation is that a reasoning process
should preserve precisely those local sections whose omission would change the
obstruction class of the evolving global solution.

\subsection{Policy Unification and Single-Signal Optimization}

Let $\tau = \{(o_t, K_t)\}_{t=1}^T$ be a trajectory. The policy gradient objective
is:
\[
  \grad_\theta\; \mathbb{E}_{\tau \sim \pi_\theta}[R(\tau)].
\]
There is no decomposition into separate objectives for reasoning and memory.
Consequently, the model learns a joint strategy in which retained state determines
future reasoning, reasoning determines reward, and reward updates both retention and
reasoning. This \emph{closes the loop} between inference and state formation---the
model becomes a self-editing dynamical system.

\subsection{Budget-Aware Interoception as Field-Awareness}

NGC prepends the eviction rate to the prompt so the model has explicit awareness of its
own resource constraint before reasoning. In RSVP terms, this is not a superficial
prompt trick but a crude form of \emph{endogenous field-awareness}: the system is told
about its own thermodynamic boundary condition before transport begins. This moves it
one step closer to a model that reasons as a function of its own internal constraint
state.

\subsection{The Replay Mask and Path-Faithful Admissibility}

The replay-mask mechanism ensures that once the model has evicted parts of its history,
training respects the exact visibility structure induced by those evictions. This is
very close to the insistence that truth is not local plausibility but fixed-point
realizability under the actual constraints of the path taken. A reasoning trace cannot
be judged by a context it never possessed. The replay mask is a path-faithful
reconstruction of the admissibility conditions under which each token was produced.

\subsection{NGC's Empirical Lyapunov Signal}

NGC's internal uncertainty functional behaves empirically as a Lyapunov function. The
paper observes that the gradient of uncertainty with respect to time is negative for
correct trajectories and positive for incorrect ones. This is not merely a heuristic
observation: it is the defining property of a Lyapunov function. We can therefore
identify
\[
  \text{uncertainty}(t) \;\approx\; \EE(x_t),
\]
and the empirical result becomes
\[
  \frac{d}{dt}\EE(x_t) \;<\; 0 \quad\text{on successful trajectories.}
\]

\begin{proposition}[Empirical Lyapunov Principle]
The internal uncertainty functional measured by NGC behaves as a Lyapunov function for
successful reasoning trajectories, providing empirical grounding for Theorem~\ref{thm:lyapunov}.
\end{proposition}

\subsection{Graph Aggregation as Computational Homotopy Colimit}

NGC's graph-based aggregation constructs a graph of reasoning paths, clusters them, and
merges them into a final structure. This is structurally:
\[
  \text{tree expansion} \to \text{selection} \to \text{collapse}.
\]
In the homotopy-theoretic picture:
\begin{itemize}[noitemsep]
  \item Different paths $\to$ different local trajectories (local sections).
  \item Clustering $\to$ identification of equivalent reasoning (homotopy identification).
  \item Summarization $\to$ quotienting (homotopy colimit).
\end{itemize}
NGC's graph reduction is therefore a computational version of taking a homotopy
colimit:
\[
  \text{stable CLIO reasoning graph} \;\equiv\; \text{bisimulation equivalence class}.
\]

\subsection{Toward Semantic Renormalization}

While NGC operates via binary retention decisions, the formalism suggests a broader
class of operators. Instead of selecting subsets $K_t$, one may consider
\emph{semantic renormalization operators}:
\[
  M_t : X_t \;\mapsto\; \tilde{X}_t
\]
that perform semantic compression, merging, or coarse-graining while preserving
invariants relevant to realizability:
\[
  \mathcal{I}(M_t(X)) \;=\; \mathcal{I}(X).
\]
The generalized CLIO operator is:
\[
  \CC_M(X) \;=\; \mathbb{E}_{o \sim \pi_\theta(\cdot \mid M(X))}[\mathcal{F}(M(X), o)],
\]
supporting compression, merging, and abstraction rather than binary deletion. This
defines a renormalization flow over semantic states, in which cognition operates by
preserving invariants while reducing representational complexity:
\[
  \text{evict} \;\subset\; \text{restrict} \;\subset\; \text{quotient} \;\subset\; \text{renormalize}.
\]
Binary garbage collection is only the crudest member of this family. A more mature
cognitive architecture would perform \emph{homotopy-aware renormalization}, preserving
invariants while collapsing redundant or obstruction-producing structure.

The deepest way to state the significance of NGC inside the present framework is this:
it is an early empirical example of a system in which cognition is being forced to
acknowledge that persistence is costly and that not all locally generated structure
deserves continued ontological status. What survives is what remains useful under
constraint. That is a weak form of the broader claim that realization is the residue
left after incompatible structure has been decimated.

\section{Worked Example: Binding and Rupture in a Finite Constraint System}

We present a minimal example illustrating the Binding Invariant Theorem concretely,
exhibiting a system with a non-trivial binding invariant, a closure event that preserves
it, and a closure event that destroys it.

\subsection{The Constraint System}

Let $T = \{t_0 \leq t_1 \leq t_2\}$. Define:
\[
  \CC_{t_0} = \{A\},\quad
  \CC_{t_1} = \{A, B\},\quad
  \CC_{t_2} = \{A, B, C\},
\]
with restriction maps $\rho_{t_1 t_0}(A,B) = A$ and $\rho_{t_2 t_1}(A,B,C) = (A,B)$.
The constraint history at $t_2$ is:
\[
  \HH_{t_2}(S) \;=\; \varprojlim\bigl\{\CC_{t_0} \leftarrow \CC_{t_1} \leftarrow \CC_{t_2}\bigr\}
  \;=\; \bigl\{(A,\; (A,B),\; (A,B,C))\bigr\}.
\]
A valid trajectory $\sigma(t_i) = \bigwedge_{j \leq i}\CC_{t_j}$ exists:
$\Gamma(T, \FF) \neq \varnothing$, hence $\BI_{t_2}(S) \neq \varnothing$.

\subsection{Preserving Closure Event}

Adjoining $D$ consistent with $(A,B,C)$: $\CC_{t_2^+} = \{A,B,C,D\}$. The binding
section extends, $\BI_{t_2^+}(S) \neq \varnothing$, and the system continues.

\subsection{Rupture Event}

Imposing $\CC_{t_2'} = \{A, B, \lnot A\}$: since $A \wedge \lnot A = \bot$, no
consistent state exists, no trajectory extends to $t_2'$, and $\BI_{t_2'}(S) =
\varnothing$. By the Rupture Criterion (Corollary~\ref{cor:rupture}), this is
replacement, not continuation.

\subsection{Partial Rupture / Ambiguity}

Imposing $\CC_{t_2''} = \{A, B, C \vee E\}$: multiple extensions exist, global
sections exist but are non-unique, and $\BI_{t_2''}(S)$ has multiple elements.
This is \emph{indeterminate continuation}: the system persists, but its future identity
is underdetermined.

\subsection{Diagrammatic Form}

The three cases correspond to the existence, non-existence, and multiplicity of
compatible cones over the inverse system.

\medskip
\noindent\textit{Binding (normal continuation).} The inverse limit
$\HH_{t_2}(S)$ is the apex of a compatible cone:
\[
\begin{tikzcd}[column sep=2.2cm, row sep=1.8cm]
  \HH_{t_2}(S)
    \arrow[d, dashed, "\pi_2"']
    \arrow[dr, dashed, "\pi_1"]
  & \\
  \{A,B,C\}
    \arrow[r, "\rho_{t_2 t_1}"']
  & \{A,B\}
    \arrow[r, "\rho_{t_1 t_0}"']
  & \{A\}
\end{tikzcd}
\]
A compatible binding section exists; $\BI_{t_2}(S) \neq \varnothing$.

\medskip
\noindent\textit{Rupture.} Imposing $\lnot A$ alongside $A$ destroys compatibility:
\[
\begin{tikzcd}[column sep=2.2cm, row sep=1.8cm]
  \{A,B,\lnot A\}
    \arrow[d, "\rho_{t_2' t_1}"']
    \arrow[dr, "\rho_{t_2' t_0}"]
  & \\
  \{A,B\}
    \arrow[r, "\rho_{t_1 t_0}"']
  & \{A\}
\end{tikzcd}
\]
No object $\HH_{t_2'}(S)$ completes this as an inverse limit: any putative section
would require $A \wedge \lnot A = \bot$. Hence $\BI_{t_2'}(S) = \varnothing$.

\medskip
\noindent\textit{Ambiguity.} The disjunctive constraint $C \vee E$ admits multiple extensions:
\[
\begin{tikzcd}[column sep=2.2cm, row sep=1.8cm]
  \{A,B,C\vee E\}
    \arrow[d, "\rho_{t_2'' t_1}"']
    \arrow[dr, "\rho_{t_2'' t_0}"]
  & \\
  \{A,B\}
    \arrow[r, "\rho_{t_1 t_0}"']
  & \{A\}
\end{tikzcd}
\]
Compatible cones exist, but not uniquely: $\BI_{t_2''}(S)$ is non-empty and
non-singleton. The system persists with underdetermined continuation.

\begin{remark}
These diagrams make visible the core distinction of the theory. Preservation
corresponds to extendability of cones. Rupture corresponds to non-existence of a
compatible cone. Ambiguity corresponds to multiplicity of cones. The binding invariant
is therefore best understood not as a substance carried through time, but as the
existence and structure of a limiting cone over the system's constraint history. This
finite example is the discrete analogue of cohomological obstruction: incompatibility
of constraints corresponds to non-vanishing obstruction classes, while compatible
extensions correspond to trivialization.
\end{remark}

\section{Open Problems}

Three problems would convert the framework's correspondences into fully tight theorems.

\textbf{Problem 1 (Lyapunov stability).} Prove that CLIO updates induce a
bisimulation-invariant Lyapunov function $V(x) = \EE(x) - \EE(x^*)$ satisfying
$V(\CC(x)) \leq V(x)$ for all $x$, with equality iff $x = x^*$, without requiring
the assumptions A2--A4 as inputs---i.e., derive them from RSVP field dynamics.

\textbf{Problem 2 (Fisher metric).} Define a natural Fisher information metric on
$\MM$ making the gradient flow of $\EE$ explicit as a natural gradient flow in the
sense of Amari. The entropy term $\mu S(x)$ suggests $\MM$ carries natural information
geometry; the Fisher metric would make the gradient flow $\partial_t x = -\grad\EE(x)$
the most efficient descent with respect to the geometry of distributions.

\textbf{Problem 3 (LCC bisimulation as homotopy---converse).} The forward direction of
Theorem~\ref{thm:bisim-homotopy} is proved. The converse---that homotopy equivalence
of descent trajectories implies LCC bisimulation---requires a realizability result
connecting the derived-stack semantics of $\MM$ to the operational semantics of LCC.

A solution to all three would unify process algebra, derived algebraic geometry, and
learning dynamics into a single formal object: the metric space of CLIO-observational
equivalence classes under Fisher-information gradient flow.

\section{Conclusion}

We have developed the binding invariant through six interlocking formalisms.

\textbf{Part~I} (sheaf and topos theory): the binding invariant is the global section
of the subobject classifier over the constraint history topos. The Binding Invariant
Theorem establishes four equivalent conditions; three corollaries characterize
uniqueness, rupture, and stability.

\textbf{Part~II} (derived algebraic geometry and variational dynamics): ambiguity is
higher homotopy in a derived semantic stack; the self is the minimum-energy homotopy
class of a multi-scale cohomological energy functional; synthetic grounding is gradient
descent; return-under-constraint is the geometric stability of the fixed point. Identity
is dynamical equilibrium, not static persistence.

\textbf{Part~III} (process algebra): LCC constraint consumption is cohomological
obstruction projection; bisimulation is gradient-flow equivalence; barbed congruence is
invariance under irreversible descent. The Triple Characterization Theorem
(Theorem~\ref{thm:char}) proves equivalence of all three levels.

\textbf{Part~IV} (monoidal category theory): CLIO is a lax monoidal endofunctor on
$\Sem$; cognitive states are fixed-point algebras of CLIO; the sheaf-theoretic lift
makes CLIO an obstruction-descending functor on $\PSh(\Omega)$.

\textbf{Part~V} (formal theorems): $\EE$ is a bisimulation-invariant Lyapunov function
(Theorem~\ref{thm:lyapunov}); LCC bisimulation corresponds to homotopy equivalence of
descent trajectories (Theorem~\ref{thm:bisim-homotopy}); quotient descent on
observational classes is simultaneously operationally meaningful and geometrically
coherent (Theorem~\ref{thm:quotient-descent}).

\textbf{Part~VI} (concrete instantiation): NGC is a concrete empirical case of CLIO
dynamics, its ``grow-then-evict'' mechanism is admissible-subspace projection in RSVP
field coordinates, its internal uncertainty is an empirical Lyapunov signal, its graph
aggregation is a computational homotopy colimit, and its limitation---binary
retention---points toward semantic renormalization as the next architecture.

The deepest result is the convergence of all six levels on the same object. The binding
invariant is not an artifact of any one formal choice but the invariant structure that
appears at every level of description of a system that undergoes irreversible change.
The framework is not speculative: it is the correct mathematical language for what NGC
is already doing. And the framework is not merely interpretive: it generates three open
problems whose resolution would close the loop between process algebra, derived geometry,
and machine learning.

\begin{quote}
\itshape
The binding invariant is the minimum-energy homotopy class that survives irreversible
projection---the observational equivalence class of a system under cohomological descent
on its constraint structure. A system is ``itself'' if it lies in a homotopy class
invariant under descent on a cohomological energy functional defined by its constraints.
\end{quote}

\vspace{2em}
\hrule
\vspace{10em}

\noindent\textbf{Acknowledgments.}
This essay develops ideas first articulated in \emph{The Compiled Self},
\emph{Return Under Constraint}, and the Yarncrawler framework. Parts~III--V draw on
Haemmerl\'e's formalization of Linear Concurrent Constraints. Part~VI engages with
Li, Hamid, Fox, and Goodman's Neural Garbage Collection.

\vspace{1em}
\noindent\textbf{Keywords.}
binding invariant, constraint closure, irreversibility, sheaf theory, topos theory,
derived stacks, $\infty$-groupoids, cohomological energy functional, gradient flow,
synthetic grounding, CLIO, RSVP field theory, Yarncrawler, linear concurrent
constraints, bisimulation, barbed congruence, observational equivalence, neural garbage
collection, semantic renormalization, monoidal endofunctor, fixed-point algebra.

\newpage
\begin{thebibliography}{99}

\bibitem{maclane}
S.~Mac Lane,
\textit{Categories for the Working Mathematician},
Springer, 2nd ed., 1998.

\bibitem{johnstone}
P.~T.~Johnstone,
\textit{Sketches of an Elephant: A Topos Theory Compendium},
Oxford University Press, 2002.

\bibitem{lurie}
J.~Lurie,
\textit{Higher Topos Theory},
Princeton University Press, 2009.

\bibitem{riehl}
E.~Riehl,
\textit{Category Theory in Context},
Dover, 2016.

\bibitem{spivak}
D.~I.~Spivak,
\textit{Category Theory for the Sciences},
MIT Press, 2014.

\bibitem{girard}
J.-Y.~Girard,
Linear Logic,
\textit{Theoretical Computer Science}, 50 (1987), 1--102.

\bibitem{saraswat}
V.~A.~Saraswat,
\textit{Concurrent Constraint Programming},
MIT Press, 1993.

\bibitem{saraswatlincoln}
V.~A.~Saraswat and P.~Lincoln,
Higher-Order, Linear, Concurrent Constraint Programming,
Technical Report, Xerox PARC, 1992.

\bibitem{haemmerle2011}
R.~Haemmerl\'e,
Observational Equivalences for Linear Logic Concurrent Constraint Languages,
\textit{Theory and Practice of Logic Programming},
11(4--5): 469--485, 2011.
Also available as arXiv:1108.0329.

\bibitem{milner}
R.~Milner,
\textit{Communication and Concurrency},
Prentice Hall, 1989.

\bibitem{parrowvictor}
J.~Parrow and B.~Victor,
The Fusion Calculus: Expressiveness and Symmetry in Mobile Processes,
\textit{Proceedings of LICS}, 1998.

\bibitem{landauer}
R.~Landauer,
Irreversibility and Heat Generation in the Computing Process,
\textit{IBM Journal of Research and Development}, 1961.

\bibitem{bennett}
C.~H.~Bennett,
The Thermodynamics of Computation,
\textit{International Journal of Theoretical Physics}, 1982.

\bibitem{lloyd}
S.~Lloyd,
Ultimate Physical Limits to Computation,
\textit{Nature}, 406 (2000), 1047--1054.

\bibitem{jacobson}
T.~Jacobson,
Thermodynamics of Spacetime: The Einstein Equation of State,
\textit{Physical Review Letters}, 75 (1995), 1260--1263.

\bibitem{verlinde}
E.~Verlinde,
On the Origin of Gravity and the Laws of Newton,
\textit{Journal of High Energy Physics}, 2011.

\bibitem{ngc2026}
M.~Y.~Li, J.~I.~Hamid, E.~B.~Fox, and N.~D.~Goodman,
Neural Garbage Collection: Learning to Forget while Learning to Reason,
arXiv:2604.18002, 2026.

\bibitem{abramsky}
S.~Abramsky and N.~Tzevelekos,
Introduction to Categories and Categorical Logic,
in \textit{New Structures for Physics}, Springer, 2011.

\bibitem{sutton2019}
R.~S.~Sutton,
The Bitter Lesson, 2019.

\end{thebibliography}

\end{document}
